% Section 7 of the manuscript: per-text agreement (question 3). Every result value is read with \rv{} from
% generated/values.tex or given in a generated table (tools/results_tables.py).

\section{Results: per-text agreement}
\label{sec:results-text}

Table~\ref{tab:text-agreement} and Figure~\ref{fig:text-agreement} give, per cell and speed, the agreement of the four
per-text measures in the primary scope: the bias and the limits of agreement of the difference between the logger and
the reference over texts of eight trials, and ICC(A,1), each with an interval that resamples blocks. These values hold
for this design's numbers of trials per text and are not compared between speeds; the ICC relates the variation between
texts to the total variation, so it also depends on how much the texts of this design differ, and it is read together
with the limits of agreement and with the reference's standard deviation between texts, which the table gives. The
design gives most trials of the primary scope one session, so the reference's conversions per trial were the same in
every text with a value on both sides in \rv{q3.conversions_per_trial.constant.strata} of the
\rv{q3.conversions_per_trial.constant.of} strata of speed and cell; there the ICC is zero whatever the logger writes and
is reported without a value, next to the registered value (deviation D11, made after the results had been seen), and
these strata are read by their limits of agreement only. In \rv{q3.conversions_per_trial.onetext.strata} further strata,
the middle and the fast speed with live conversion on, the reference differed from one conversion per trial in a single
text, so their ICCs are near zero for the same reason. Where the reference's values vary little between texts, as for
these strata and for the fraction of time converting on macOS, a low ICC also reflects the small differences between the
texts of this design; the size of the disagreement is read from the limits of agreement, which the table sets beside the
reference's standard deviation between texts.

The bias of the mean conversion time was \rv{macos-standard-live-off.l2.mean_dwell_ms.slow.bias},
\rv{macos-standard-live-off.l2.mean_dwell_ms.mid.bias} and \rv{macos-standard-live-off.l2.mean_dwell_ms.fast.bias}
(slow, middle and fast speed) with live conversion off, \rv{macos-standard-live-on.l2.mean_dwell_ms.slow.bias},
\rv{macos-standard-live-on.l2.mean_dwell_ms.mid.bias} and \rv{macos-standard-live-on.l2.mean_dwell_ms.fast.bias} with
it on, and \rv{windows-msime-tsf.l2.mean_dwell_ms.slow.bias}, \rv{windows-msime-tsf.l2.mean_dwell_ms.mid.bias} and
\rv{windows-msime-tsf.l2.mean_dwell_ms.fast.bias} on Windows. The bias of the abandonment rate and that of the
fraction of time converting were negative, with intervals below zero, at every speed in every cell: the logger's texts
had a lower abandonment rate and a lower fraction of time converting than the reference's. The bias of the conversions
per trial differed in sign between the speeds on macOS and was positive at every speed on Windows
(Table~\ref{tab:text-agreement}).

The supplement gives the same measures over all trials, the values pooled over the speeds (for comparison only, because
a difference in level between the speeds inflates a pooled ICC), the bias under the three mixtures of speeds of the
registration, and the per-trial differences per scenario and speed with the mean on each side, so that the scenarios
can be weighted differently. The views of the per-trial differences that leave out the trials of the carry layers need
the prediction and are given for the macOS cells only.

% Generated by manuscript/tools/results_tables.py from the result files; do not edit by hand.
\begin{sidewaystable}[p]
{\footnotesize
\caption{Per-text agreement per cell, measure and speed, primary scope: bias and limits of agreement of the per-text difference (observation minus reference) and ICC(A,1), each with a 95\% interval that resamples blocks; Texts is the number of texts with a value on both sides, and Reference SD the standard deviation of the reference's values between those texts. An estimate printed without an interval has its interval withheld (the reason is in the result file); --: no value; -- (reference constant; registered \ldots): no ICC(A,1), because the reference's value is the same in every text, so the ICC is zero whatever the logger writes, or undefined if the logger writes that same value in every text too (section 6.11 of the registration; deviation D11); the value that the frozen code returns, the registered result, follows \emph{registered}. The values hold for this design's numbers of trials per text and are not compared between speeds}
\label{tab:text-agreement}
\begin{tabular}{L{0.13}L{0.13}L{0.06}L{0.06}L{0.09}L{0.13}L{0.13}L{0.13}L{0.13}}
\toprule
Cell & Measure & Speed & Texts & Reference SD & Bias & Lower limit & Upper limit & ICC(A,1) \\
\midrule
macOS, live conversion off & Conversions per trial & slow & 117 & 0.000 & 0.034 [0.014, 0.057] & $-$0.179 [$-$0.221, $-$0.136] & 0.248 [0.180, 0.309] & -- (reference constant; registered 0.000 [0.000, 0.000]) \\
 &  & mid & 117 & 0.000 & $-$0.003 [$-$0.018, 0.009] & $-$0.181 [$-$0.220, $-$0.137] & 0.175 [0.134, 0.215] & -- (reference constant; registered 0.000 [0.000, 0.000]) \\
 &  & fast & 117 & 0.000 & $-$0.051 [$-$0.068, $-$0.032] & $-$0.286 [$-$0.318, $-$0.248] & 0.184 [0.150, 0.215] & -- (reference constant; registered 0.000 [0.000, 0.000]) \\
 & Mean conversion time (ms) & slow & 117 & 2,632 & $-$971 [$-$1,262, $-$691] & $-$3,841 [$-$4,611, $-$3,042] & 1,900 [1,358, 2,423] & 0.809 [0.700, 0.888] \\
 &  & mid & 117 & 1,891 & $-$379 [$-$558, $-$179] & $-$2,663 [$-$3,074, $-$2,241] & 1,905 [1,414, 2,390] & 0.811 [0.712, 0.877] \\
 &  & fast & 117 & 1,395 & 86 [$-$95, 258] & $-$2,025 [$-$2,338, $-$1,660] & 2,197 [1,846, 2,515] & 0.714 [0.604, 0.792] \\
 & Abandonment rate & slow & 117 & 0.203 & $-$0.012 [$-$0.018, $-$0.006] & $-$0.078 [$-$0.104, $-$0.047] & 0.054 [0.034, 0.070] & 0.984 [0.966, 0.995] \\
 &  & mid & 117 & 0.149 & $-$0.034 [$-$0.047, $-$0.022] & $-$0.178 [$-$0.239, $-$0.121] & 0.110 [0.074, 0.150] & 0.848 [0.725, 0.928] \\
 &  & fast & 117 & 0.162 & $-$0.068 [$-$0.081, $-$0.055] & $-$0.230 [$-$0.267, $-$0.194] & 0.095 [0.074, 0.114] & 0.782 [0.682, 0.851] \\
 & Fraction of time converting & slow & 117 & 0.007 & $-$0.056 [$-$0.073, $-$0.042] & $-$0.226 [$-$0.279, $-$0.168] & 0.113 [0.080, 0.142] & 0.002 [$-$0.011, 0.020] \\
 &  & mid & 117 & 0.004 & $-$0.051 [$-$0.065, $-$0.038] & $-$0.209 [$-$0.248, $-$0.167] & 0.107 [0.082, 0.130] & 0.003 [$-$0.009, 0.017] \\
 &  & fast & 117 & 0.012 & $-$0.043 [$-$0.060, $-$0.027] & $-$0.239 [$-$0.285, $-$0.187] & 0.153 [0.125, 0.182] & 0.000 [$-$0.017, 0.017] \\
macOS, live conversion on & Conversions per trial & slow & 82 & 0.000 & 0.050 [0.025, 0.076] & $-$0.214 [$-$0.262, $-$0.154] & 0.314 [0.215, 0.399] & -- (reference constant; registered 0.000 [0.000, 0.000]) \\
 &  & mid & 99 & 0.017 & 0.012 [$-$0.008, 0.036] & $-$0.181 [$-$0.218, $-$0.143] & 0.206 [0.162, 0.245] & 0.050 [0.000, 0.169] \\
 &  & fast & 114 & 0.012 & $-$0.026 [$-$0.047, $-$0.006] & $-$0.262 [$-$0.311, $-$0.214] & 0.210 [0.160, 0.260] & 0.004 [0.000, 0.015] \\
 & Mean conversion time (ms) & slow & 82 & 2,481 & $-$1,014 [$-$1,287, $-$770] & $-$3,553 [$-$4,316, $-$2,726] & 1,524 [1,129, 1,824] & 0.820 [0.689, 0.904] \\
 &  & mid & 99 & 1,899 & $-$173 [$-$393, 46] & $-$2,255 [$-$2,690, $-$1,764] & 1,908 [1,222, 2,540] & 0.874 [0.800, 0.929] \\
 &  & fast & 114 & 1,307 & $-$82 [$-$296, 157] & $-$2,385 [$-$2,895, $-$1,883] & 2,221 [1,548, 3,036] & 0.683 [0.532, 0.798] \\
 & Abandonment rate & slow & 82 & 0.204 & $-$0.010 [$-$0.016, $-$0.004] & $-$0.074 [$-$0.097, $-$0.046] & 0.055 [0.033, 0.070] & 0.986 [0.969, 0.995] \\
 &  & mid & 99 & 0.167 & $-$0.026 [$-$0.037, $-$0.016] & $-$0.138 [$-$0.170, $-$0.100] & 0.086 [0.066, 0.100] & 0.930 [0.874, 0.967] \\
 &  & fast & 114 & 0.159 & $-$0.058 [$-$0.073, $-$0.045] & $-$0.218 [$-$0.257, $-$0.180] & 0.101 [0.079, 0.125] & 0.805 [0.705, 0.876] \\
 & Fraction of time converting & slow & 82 & 0.009 & $-$0.047 [$-$0.058, $-$0.037] & $-$0.155 [$-$0.203, $-$0.099] & 0.062 [0.024, 0.090] & 0.023 [$-$0.006, 0.064] \\
 &  & mid & 99 & 0.007 & $-$0.024 [$-$0.031, $-$0.017] & $-$0.115 [$-$0.149, $-$0.077] & 0.067 [0.030, 0.101] & 0.037 [$-$0.008, 0.101] \\
 &  & fast & 114 & 0.013 & $-$0.053 [$-$0.074, $-$0.033] & $-$0.276 [$-$0.324, $-$0.222] & 0.170 [0.145, 0.192] & $-$0.013 [$-$0.054, 0.027] \\
Windows & Conversions per trial & slow & 117 & 0.000 & 0.167 [0.127, 0.209] & $-$0.250 [$-$0.310, $-$0.187] & 0.584 [0.475, 0.691] & -- (reference constant; registered 0.000 [0.000, 0.000]) \\
 &  & mid & 116 & 0.000 & 0.157 [0.134, 0.179] & $-$0.122 [$-$0.166, $-$0.080] & 0.437 [0.375, 0.495] & -- (reference constant; registered 0.000 [0.000, 0.000]) \\
 &  & fast & 117 & 0.061 & 0.038 [0.013, 0.065] & $-$0.257 [$-$0.319, $-$0.186] & 0.334 [0.225, 0.431] & 0.267 [0.117, 0.461] \\
 & Mean conversion time (ms) & slow & 117 & 2,460 & $-$2,201 [$-$2,579, $-$1,852] & $-$5,878 [$-$6,749, $-$5,032] & 1,476 [984, 1,976] & 0.503 [0.367, 0.622] \\
 &  & mid & 116 & 1,963 & $-$2,343 [$-$2,572, $-$2,116] & $-$5,058 [$-$5,490, $-$4,575] & 371 [$-$26, 709] & 0.361 [0.262, 0.451] \\
 &  & fast & 117 & 1,410 & $-$1,395 [$-$1,601, $-$1,207] & $-$3,604 [$-$3,979, $-$3,199] & 814 [565, 1,043] & 0.371 [0.218, 0.516] \\
 & Abandonment rate & slow & 117 & 0.212 & $-$0.074 [$-$0.097, $-$0.052] & $-$0.318 [$-$0.404, $-$0.231] & 0.171 [0.118, 0.222] & 0.708 [0.588, 0.820] \\
 &  & mid & 116 & 0.156 & $-$0.150 [$-$0.172, $-$0.127] & $-$0.398 [$-$0.435, $-$0.357] & 0.099 [0.069, 0.128] & 0.314 [0.217, 0.409] \\
 &  & fast & 117 & 0.166 & $-$0.171 [$-$0.193, $-$0.149] & $-$0.436 [$-$0.478, $-$0.389] & 0.094 [0.066, 0.118] & 0.291 [0.207, 0.372] \\
 & Fraction of time converting & slow & 117 & 0.027 & $-$0.066 [$-$0.076, $-$0.056] & $-$0.186 [$-$0.223, $-$0.145] & 0.055 [0.028, 0.076] & 0.112 [0.032, 0.242] \\
 &  & mid & 116 & 0.039 & $-$0.136 [$-$0.156, $-$0.117] & $-$0.320 [$-$0.369, $-$0.270] & 0.048 [0.020, 0.072] & 0.087 [$-$0.001, 0.175] \\
 &  & fast & 117 & 0.070 & $-$0.210 [$-$0.238, $-$0.182] & $-$0.515 [$-$0.556, $-$0.468] & 0.095 [0.054, 0.130] & 0.010 [$-$0.059, 0.079] \\
\bottomrule
\end{tabular}
}
\end{sidewaystable}

% Generated by manuscript/tools/results_tables.py from the result files; do not edit by hand.
\begin{figure}[p]
\centering
\pgfplotslegendfromname{fig:text-agreement:legend}\par\medskip
\begin{tikzpicture}
\begin{groupplot}[group style={group size=1 by 4, vertical sep=1.4cm}, fig panel, fig zero, scale only axis, width=11cm, height=2.5cm, xtick={1,2,3}, xticklabels={{macOS, live conversion off},{macOS, live conversion on},{Windows}}, xmin=0.5, xmax=3.5]
\nextgroupplot[title={Conversions per trial}, legend to name=fig:text-agreement:legend, legend columns=4]
\addlegendimage{fig slow, fig point}
\addlegendentry{slow}
\addlegendimage{fig mid, fig point}
\addlegendentry{mid}
\addlegendimage{fig fast, fig point}
\addlegendentry{fast}
\addlegendimage{black, fig limit}
\addlegendentry{limit of agreement}
\addplot[forget plot, fig interval, fig slow] table[row sep=\\] {x y\\0.75 0.0137968\\0.75 0.0565787\\0.75 nan\\0.75 -0.220809\\0.75 -0.136131\\0.75 nan\\0.75 0.179633\\0.75 0.308846\\0.75 nan\\1.75 0.0248955\\1.75 0.0760574\\1.75 nan\\1.75 -0.261969\\1.75 -0.15442\\1.75 nan\\1.75 0.215324\\1.75 0.399236\\1.75 nan\\2.75 0.127467\\2.75 0.209161\\2.75 nan\\2.75 -0.310034\\2.75 -0.187389\\2.75 nan\\2.75 0.474696\\2.75 0.690745\\2.75 nan\\};
\addplot[forget plot, fig slow, fig limit] table[row sep=\\] {x y\\0.75 -0.178767\\1.75 -0.214465\\2.75 -0.250327\\0.75 0.247509\\1.75 0.313681\\2.75 0.583864\\};
\addplot[forget plot, fig slow, fig point] table[row sep=\\] {x y\\0.75 0.0343712\\1.75 0.049608\\2.75 0.166768\\};
\addplot[forget plot, fig interval, fig mid] table[row sep=\\] {x y\\1 -0.0178577\\1 0.00922975\\1 nan\\1 -0.220274\\1 -0.137461\\1 nan\\1 0.133838\\1 0.215377\\1 nan\\2 -0.00826278\\2 0.0355079\\2 nan\\2 -0.217972\\2 -0.142543\\2 nan\\2 0.161662\\2 0.245401\\2 nan\\3 0.134052\\3 0.179142\\3 nan\\3 -0.166346\\3 -0.0797757\\3 nan\\3 0.375232\\3 0.494742\\3 nan\\};
\addplot[forget plot, fig mid, fig limit] table[row sep=\\] {x y\\1 -0.181472\\2 -0.181242\\3 -0.12198\\1 0.174879\\2 0.206037\\3 0.436799\\};
\addplot[forget plot, fig mid, fig point] table[row sep=\\] {x y\\1 -0.0032967\\2 0.0123978\\3 0.15741\\};
\addplot[forget plot, fig interval, fig fast] table[row sep=\\] {x y\\1.25 -0.0684778\\1.25 -0.0321009\\1.25 nan\\1.25 -0.318354\\1.25 -0.247791\\1.25 nan\\1.25 0.149591\\1.25 0.214642\\1.25 nan\\2.25 -0.0471251\\2.25 -0.00645323\\2.25 nan\\2.25 -0.311293\\2.25 -0.214478\\2.25 nan\\2.25 0.159809\\2.25 0.259746\\2.25 nan\\3.25 0.0132262\\3.25 0.0645655\\3.25 nan\\3.25 -0.318555\\3.25 -0.185716\\3.25 nan\\3.25 0.225438\\3.25 0.43123\\3.25 nan\\};
\addplot[forget plot, fig fast, fig limit] table[row sep=\\] {x y\\1.25 -0.28563\\2.25 -0.262432\\3.25 -0.256792\\1.25 0.18388\\2.25 0.20955\\3.25 0.333511\\};
\addplot[forget plot, fig fast, fig point] table[row sep=\\] {x y\\1.25 -0.0508751\\2.25 -0.0264411\\3.25 0.0383598\\};
\nextgroupplot[title={Mean conversion time (ms)}]
\addplot[forget plot, fig interval, fig slow] table[row sep=\\] {x y\\0.75 -1261.86\\0.75 -691.009\\0.75 nan\\0.75 -4610.86\\0.75 -3041.94\\0.75 nan\\0.75 1357.61\\0.75 2423.44\\0.75 nan\\1.75 -1287.48\\1.75 -769.965\\1.75 nan\\1.75 -4315.51\\1.75 -2725.52\\1.75 nan\\1.75 1129.49\\1.75 1824.28\\1.75 nan\\2.75 -2579.08\\2.75 -1852.1\\2.75 nan\\2.75 -6748.57\\2.75 -5032.46\\2.75 nan\\2.75 984.201\\2.75 1976.19\\2.75 nan\\};
\addplot[forget plot, fig slow, fig limit] table[row sep=\\] {x y\\0.75 -3841.17\\1.75 -3552.51\\2.75 -5878.05\\0.75 1899.84\\1.75 1523.73\\2.75 1475.88\\};
\addplot[forget plot, fig slow, fig point] table[row sep=\\] {x y\\0.75 -970.666\\1.75 -1014.39\\2.75 -2201.08\\};
\addplot[forget plot, fig interval, fig mid] table[row sep=\\] {x y\\1 -557.669\\1 -178.966\\1 nan\\1 -3073.94\\1 -2240.67\\1 nan\\1 1413.68\\1 2389.74\\1 nan\\2 -393.016\\2 45.8076\\2 nan\\2 -2690.07\\2 -1764.42\\2 nan\\2 1222.49\\2 2540.43\\2 nan\\3 -2571.92\\3 -2116.47\\3 nan\\3 -5490.45\\3 -4574.81\\3 nan\\3 -25.5713\\3 709.108\\3 nan\\};
\addplot[forget plot, fig mid, fig limit] table[row sep=\\] {x y\\1 -2662.93\\2 -2254.78\\3 -5058.22\\1 1905.33\\2 1907.89\\3 371.489\\};
\addplot[forget plot, fig mid, fig point] table[row sep=\\] {x y\\1 -378.801\\2 -173.443\\3 -2343.37\\};
\addplot[forget plot, fig interval, fig fast] table[row sep=\\] {x y\\1.25 -94.6766\\1.25 257.511\\1.25 nan\\1.25 -2338.43\\1.25 -1660.41\\1.25 nan\\1.25 1845.51\\1.25 2514.94\\1.25 nan\\2.25 -296.116\\2.25 156.763\\2.25 nan\\2.25 -2895.36\\2.25 -1883.35\\2.25 nan\\2.25 1547.61\\2.25 3036.08\\2.25 nan\\3.25 -1600.93\\3.25 -1207.2\\3.25 nan\\3.25 -3979.15\\3.25 -3198.67\\3.25 nan\\3.25 565.278\\3.25 1042.63\\3.25 nan\\};
\addplot[forget plot, fig fast, fig limit] table[row sep=\\] {x y\\1.25 -2025.32\\2.25 -2384.61\\3.25 -3604.34\\1.25 2196.62\\2.25 2221.25\\3.25 814.444\\};
\addplot[forget plot, fig fast, fig point] table[row sep=\\] {x y\\1.25 85.6454\\2.25 -81.6779\\3.25 -1394.95\\};
\nextgroupplot[title={Abandonment rate}]
\addplot[forget plot, fig interval, fig slow] table[row sep=\\] {x y\\0.75 -0.0183359\\0.75 -0.00550392\\0.75 nan\\0.75 -0.103954\\0.75 -0.0466396\\0.75 nan\\0.75 0.0339079\\0.75 0.0695433\\0.75 nan\\1.75 -0.0162794\\1.75 -0.00358063\\1.75 nan\\1.75 -0.0974907\\1.75 -0.0464943\\1.75 nan\\1.75 0.0334953\\1.75 0.0702787\\1.75 nan\\2.75 -0.0974065\\2.75 -0.0521038\\2.75 nan\\2.75 -0.404381\\2.75 -0.231365\\2.75 nan\\2.75 0.118328\\2.75 0.222407\\2.75 nan\\};
\addplot[forget plot, fig slow, fig limit] table[row sep=\\] {x y\\0.75 -0.0775651\\1.75 -0.0740273\\2.75 -0.318412\\0.75 0.0542846\\1.75 0.0546118\\2.75 0.171065\\};
\addplot[forget plot, fig slow, fig point] table[row sep=\\] {x y\\0.75 -0.0116402\\1.75 -0.0097077\\2.75 -0.0736739\\};
\addplot[forget plot, fig interval, fig mid] table[row sep=\\] {x y\\1 -0.0468264\\1 -0.0222586\\1 nan\\1 -0.238715\\1 -0.120698\\1 nan\\1 0.0736147\\1 0.149951\\1 nan\\2 -0.0372833\\2 -0.0155313\\2 nan\\2 -0.169692\\2 -0.100276\\2 nan\\2 0.0664793\\2 0.0997035\\2 nan\\3 -0.171625\\3 -0.127335\\3 nan\\3 -0.435022\\3 -0.356656\\3 nan\\3 0.0691825\\3 0.127661\\3 nan\\};
\addplot[forget plot, fig mid, fig limit] table[row sep=\\] {x y\\1 -0.178358\\2 -0.13783\\3 -0.398436\\1 0.110083\\2 0.0856654\\3 0.0987183\\};
\addplot[forget plot, fig mid, fig point] table[row sep=\\] {x y\\1 -0.0341372\\2 -0.0260823\\3 -0.149859\\};
\addplot[forget plot, fig interval, fig fast] table[row sep=\\] {x y\\1.25 -0.0808639\\1.25 -0.0554578\\1.25 nan\\1.25 -0.266532\\1.25 -0.1935\\1.25 nan\\1.25 0.0736502\\1.25 0.113685\\1.25 nan\\2.25 -0.0728879\\2.25 -0.0449424\\2.25 nan\\2.25 -0.256723\\2.25 -0.180276\\2.25 nan\\2.25 0.0787637\\2.25 0.124625\\2.25 nan\\3.25 -0.192895\\3.25 -0.148815\\3.25 nan\\3.25 -0.477632\\3.25 -0.389368\\3.25 nan\\3.25 0.0663391\\3.25 0.118487\\3.25 nan\\};
\addplot[forget plot, fig fast, fig limit] table[row sep=\\] {x y\\1.25 -0.230252\\2.25 -0.217832\\3.25 -0.436111\\1.25 0.0948428\\2.25 0.101054\\3.25 0.0937557\\};
\addplot[forget plot, fig fast, fig point] table[row sep=\\] {x y\\1.25 -0.0677045\\2.25 -0.058389\\3.25 -0.171178\\};
\nextgroupplot[title={Fraction of time converting}]
\addplot[forget plot, fig interval, fig slow] table[row sep=\\] {x y\\0.75 -0.0727727\\0.75 -0.0419352\\0.75 nan\\0.75 -0.278801\\0.75 -0.16807\\0.75 nan\\0.75 0.0801202\\0.75 0.14193\\0.75 nan\\1.75 -0.0583606\\1.75 -0.0365913\\1.75 nan\\1.75 -0.203414\\1.75 -0.0991124\\1.75 nan\\1.75 0.0236671\\1.75 0.0902096\\1.75 nan\\2.75 -0.0757142\\2.75 -0.0555747\\2.75 nan\\2.75 -0.223332\\2.75 -0.144563\\2.75 nan\\2.75 0.0279813\\2.75 0.0763184\\2.75 nan\\};
\addplot[forget plot, fig slow, fig limit] table[row sep=\\] {x y\\0.75 -0.226043\\1.75 -0.15533\\2.75 -0.18638\\0.75 0.113162\\1.75 0.062031\\2.75 0.0552155\\};
\addplot[forget plot, fig slow, fig point] table[row sep=\\] {x y\\0.75 -0.0564405\\1.75 -0.0466495\\2.75 -0.0655823\\};
\addplot[forget plot, fig interval, fig mid] table[row sep=\\] {x y\\1 -0.0647352\\1 -0.037892\\1 nan\\1 -0.24751\\1 -0.166924\\1 nan\\1 0.0818906\\1 0.129647\\1 nan\\2 -0.0310215\\2 -0.0169641\\2 nan\\2 -0.148586\\2 -0.0769847\\2 nan\\2 0.0300855\\2 0.101428\\2 nan\\3 -0.156174\\3 -0.116918\\3 nan\\3 -0.369001\\3 -0.270129\\3 nan\\3 0.0200337\\3 0.0721026\\3 nan\\};
\addplot[forget plot, fig mid, fig limit] table[row sep=\\] {x y\\1 -0.209188\\2 -0.114821\\3 -0.320378\\1 0.107008\\2 0.0667371\\3 0.0483024\\};
\addplot[forget plot, fig mid, fig point] table[row sep=\\] {x y\\1 -0.0510901\\2 -0.0240419\\3 -0.136038\\};
\addplot[forget plot, fig interval, fig fast] table[row sep=\\] {x y\\1.25 -0.0598186\\1.25 -0.0271184\\1.25 nan\\1.25 -0.28509\\1.25 -0.187262\\1.25 nan\\1.25 0.124557\\1.25 0.181799\\1.25 nan\\2.25 -0.074464\\2.25 -0.0333915\\2.25 nan\\2.25 -0.323565\\2.25 -0.22158\\2.25 nan\\2.25 0.145255\\2.25 0.191593\\2.25 nan\\3.25 -0.238401\\3.25 -0.181827\\3.25 nan\\3.25 -0.556206\\3.25 -0.468247\\3.25 nan\\3.25 0.0544072\\3.25 0.129845\\3.25 nan\\};
\addplot[forget plot, fig fast, fig limit] table[row sep=\\] {x y\\1.25 -0.23921\\2.25 -0.275999\\3.25 -0.514636\\1.25 0.153092\\2.25 0.170183\\3.25 0.0950716\\};
\addplot[forget plot, fig fast, fig point] table[row sep=\\] {x y\\1.25 -0.043059\\2.25 -0.0529079\\3.25 -0.209782\\};
\end{groupplot}
\end{tikzpicture}
\caption{Per-text agreement of the four measures per cell and speed in the primary scope: the bias (marks) and the lower and upper limits of agreement (dashes) of the per-text difference (observation minus reference), each with a 95\% interval that resamples blocks. The values hold for this design's numbers of trials per text and are not compared between speeds, and the values of the Windows cell estimate the agreement with a different construct and are not compared with those of the macOS cells; a measure with no text in the scope is not drawn. Table~\ref{tab:text-agreement} gives the values with ICC(A,1). The figure is drawn from the result files by code written and reviewed with an AI tool (Claude) and checked by its tests; see the AI declaration}
\label{fig:text-agreement}
\end{figure}

